\chapter{Conclusioni}

KompozeR e' stato progettato come una web application modulare in cui il valore principale risiede nel configuratore CAD e nella sua integrazione con catalogo, carrello, chat, notifiche e amministrazione. L'analisi del repository reale mostra un progetto piu' concreto di quello abbozzato nelle note iniziali: il frontend e' una SPA Vue 3 con Pinia, il gateway centralizza autenticazione e proxy, i servizi backend sono separati per responsabilita' e il CAD e' un core fortemente orientato alla consistenza del modello.

Il punto piu' importante emerso dal doppio controllo tra note e codice e' che la documentazione deve descrivere il progetto implementato, non quello solo immaginato nelle prime bozze. Per questo il documento insiste sui path realmente esposti, sulla logica step-based del CAD, sui payload effettivi, sulla presenza di un BFF nel gateway e sulla suite di test disponibile nel repository.

In sintesi, il progetto presenta una struttura coerente per un'applicazione web moderna e un buon grado di modularita'. Gli sviluppi futuri piu' naturali riguardano il rafforzamento dei test UI, l'eventuale estensione del CAD collaborativo e un consolidamento ulteriore delle verifiche end-to-end di accessibilita' e responsive design.
